\subsection{Выбор программных средств для серверной части и веб-панели}

В качестве СУБД выбрана PostgreSQL с расширением PostGIS~--- оно позволяет выполнять геопространственные запросы непосредственно на уровне базы данных. Функция ST\_DWithin расширения PostGIS позволяет эффективно искать обращения в заданном радиусе от указанной точки, что является основой алгоритма дедупликации.

Серверная часть реализована на C\# с использованием ASP.NET Core и Entity Framework Core. Подход Code-First позволяет описывать модель данных в виде классов C\# и автоматически генерировать миграции. Выбор обусловлен следующими преимуществами: кроссплатформенность среды выполнения .NET, встроенная поддержка HTTPS и JWT-аутентификации, интеграция с OpenAPI (Swagger) для автодокументирования API, производительность и масштабируемость.

Для веб-панели диспетчера выбраны React и TypeScript с библиотекой Яндекс.Карты (JavaScript API) для отображения интерактивной карты. Сборка выполняется инструментом Vite. Выбор React обусловлен компонентным подходом, TypeScript обеспечивает статическую типизацию и раннее обнаружение ошибок, Яндекс.Карты предоставляют готовый картографический сервис с актуальными данными по территории России.

\subsection{Реализация базы данных и серверной логики}

Работа с базой данных выполнена с использованием подхода Code-First в Entity Framework Core. Реализованы классы-сущности, конфигурации через Fluent API и контекст базы данных SmartCityDbContext. Конфигурации устанавливают ограничения длины строковых полей, уникальные индексы, правила каскадного удаления и составные ограничения.

Для инкапсуляции логики доступа к данным реализован паттерн <<Репозиторий>>. Специализированный ReportRepository включает методы фильтрации по статусу, получения обращения со связанными данными (категория, медиафайлы, история статусов) и поиска дубликатов по координатам. Репозитории объединены через UnitOfWork~--- класс, обеспечивающий атомарность транзакций и единый контекст базы данных для всех репозиториев в рамках одного запроса.

Реализованы контроллеры программного интерфейса:
\begin{itemize}
    \item ReportsController~--- CRUD-операции с обращениями, поиск дубликатов по координатам и категории, подтверждение обращения другим пользователем, получение списка обращений с фильтрацией по статусу;
    \item CategoriesController~--- получение, создание, редактирование и удаление категорий проблем;
    \item ServicesController~--- справочник муниципальных служб, привязка служб к категориям;
    \item AuthController~--- регистрация, аутентификация по логину/паролю, аутентификация через Google, получение текущего пользователя, восстановление пароля;
    \item EventsController~--- городские события с модерацией.
\end{itemize}

\subsection{Ключевые алгоритмы серверной части}

\subsubsection{Алгоритм дедупликации обращений}

При создании нового обращения среди активных обращений производится поиск записей с совпадающей категорией в радиусе 50~м от указанной точки. Геопространственный поиск выполняется функцией ST\_DWithin расширения PostGIS. При обнаружении дубликата пользователю предлагается подтвердить существующее обращение вместо создания нового. На рисунке~\ref{fig:alg-dedup} представлена блок-схема алгоритма.

\begin{figure}[H]
    \centering
    \includegraphics[width=0.5\textwidth]{alg-dedup.png}
    \caption{Алгоритм дедупликации обращений}
    \label{fig:alg-dedup}
\end{figure}


\subsubsection{Алгоритм обработки смены статуса}

Жизненный цикл обращения включает пять статусов: New, Accepted, InProgress, Resolved, Rejected. При каждом изменении статуса создаётся запись в ReportStatusHistory с указанием даты, времени, идентификатора оператора и комментария. Система выполняет валидацию допустимости перехода между статусами. На рисунке~\ref{fig:alg-status} представлена блок-схема.

\begin{figure}[H]
    \centering
    \includegraphics[width=0.5\textwidth]{alg-status.png}
    \caption{Алгоритм обработки смены статуса обращения}
    \label{fig:alg-status}
\end{figure}

\subsection{Реализация веб-панели диспетчера}

Веб-панель построена на React с TypeScript и сборщиком Vite. Управление состоянием реализовано через React Context. Интерфейс включает несколько основных страниц, описанных ниже.

Дашборд (рисунок~\ref{fig:web-dashboard}) отображает сводную статистику по обращениям: общее количество, количество по каждому статусу, графики динамики поступления и распределения по категориям. Для визуализации статистики используется библиотека Recharts.

\begin{figure}[H]
    \centering
    \setlength{\fboxrule}{0.4pt}
    \setlength{\fboxsep}{0pt}
    \fbox{\includegraphics[width=0.96\linewidth,height=0.42\textheight,keepaspectratio]{impl-dashboard.png}}
    \caption{Дашборд веб-панели диспетчера}
    \label{fig:web-dashboard}
\end{figure}

Интерактивная карта обращений (рисунок~\ref{fig:web-map}) отображает маркеры обращений, цвет которых определяется статусом (новое~--- синий, принято~--- оранжевый, в работе~--- жёлтый, решено~--- зелёный). Реализованы фильтрация по статусу и категории, кластеризация маркеров при уменьшении масштаба и всплывающие окна с краткой информацией об обращении.

\begin{figure}[H]
    \centering
    \setlength{\fboxrule}{0.4pt}
    \setlength{\fboxsep}{0pt}
    \fbox{\includegraphics[width=0.96\linewidth,height=0.42\textheight,keepaspectratio]{impl-map.png}}
    \caption{Интерактивная карта обращений}
    \label{fig:web-map}
\end{figure}

Страница обращений (рисунок~\ref{fig:web-reports}) содержит табличный список обращений с вкладками по статусам (<<Все>>, <<Новые>>, <<Принятые>>, <<В работе>>, <<Решённые>>, <<Отклонённые>>), поиском по описанию и переходом к карточке обращения.

\begin{figure}[H]
    \centering
    \setlength{\fboxrule}{0.4pt}
    \setlength{\fboxsep}{0pt}
    \fbox{\includegraphics[width=0.96\linewidth,height=0.42\textheight,keepaspectratio]{impl-reports.png}}
    \caption{Страница обращений веб-панели}
    \label{fig:web-reports}
\end{figure}

Карточка обращения (рисунок~\ref{fig:web-report-card}) отображает полную информацию об обращении: фотографию, описание, категорию, мини-карту с маркером, историю статусов и контактные данные ответственной службы. Диспетчер может изменить статус обращения через форму на этой странице.

\begin{figure}[H]
    \centering
    \setlength{\fboxrule}{0.4pt}
    \setlength{\fboxsep}{0pt}
    \fbox{\includegraphics[width=0.96\linewidth,height=0.42\textheight,keepaspectratio]{impl-report-card.png}}
    \caption{Карточка обращения}
    \label{fig:web-report-card}
\end{figure}

Справочник категорий проблем (рисунок~\ref{fig:web-categories}) обеспечивает просмотр, добавление, редактирование и удаление категорий. Каждая категория содержит наименование, код и связь с ответственными службами.

\begin{figure}[H]
    \centering
    \setlength{\fboxrule}{0.4pt}
    \setlength{\fboxsep}{0pt}
    \fbox{\includegraphics[width=0.96\linewidth,height=0.42\textheight,keepaspectratio]{impl-categories.png}}
    \caption{Справочник категорий проблем}
    \label{fig:web-categories}
\end{figure}

Страница муниципальных служб (рисунок~\ref{fig:web-services}) содержит перечень городских служб с контактной информацией (телефон, адрес, режим работы) и привязкой к категориям проблем.

\begin{figure}[H]
    \centering
    \setlength{\fboxrule}{0.4pt}
    \setlength{\fboxsep}{0pt}
    \fbox{\includegraphics[width=0.96\linewidth,height=0.42\textheight,keepaspectratio]{impl-services.png}}
    \caption{Страница муниципальных служб}
    \label{fig:web-services}
\end{figure}
